\subsection{BASES OF AN ALGEBRA. MULTIPLICATION TABLE}

By definition, a basis of an A-algebra E is a basis of E for its A-module structure. Let \((a_{i})_{i\in I}\) be a basis of E; there exists a unique family \((\gamma_{ij}^{k})_{(i,j,k)\in I\times I\times I}\) of elements of the ring A such that for every ordered pair \((i,j)\in I\times I\) , the set of \(k\in I\) such that \(\gamma_{ij}^{k}\neq 0\) is finite and

\[
a _ {i} a _ {j} = \sum_ {k \in \mathbf {I}} \gamma_ {i j} ^ {k} a _ {k}.\tag{7}
\]

The \(\gamma_{ij}^{k}\) are called the constants of structure of the algebra E with respect to the basis \((a_{i})\) and relations (7) constitute the multiplication table of the algebra E (relative to the basis \((a_{i})\) ).

Relations (7) can be imagined written down by setting out the right hand sides of these relations in a square table

\(\cdots\)

\(a_j\)

\(\cdots\)

\(\vdots\)

\(a_j\)

\(\sum_{k} \gamma_{ij}^{k} a_k\)

\(\vdots\)

it being understood that the element appearing in the row of index i and the column of index j is equal to the product \(a_{i}a_{j}\) .

Conversely, given an A-module E, a basis \((a_{i})_{i\in I}\) of E and a family \((\gamma_{ij}^{k})\) of elements of A such that for every ordered pair \((i,j)\in I\times I\) the set of \(k\in I\) such that \(\gamma_{ij}^{k}\neq 0\) is finite, then there is on E one and only one A-algebra structure under which relations (7) hold, since the A-module \(\mathbf{E}\otimes_{\mathbf{A}}\mathbf{E}\) is free and admits as basis \((a_{i}\otimes a_{j})_{(i,j)\in I\times I}\) (cf.~II, § 3, no. 7, Corollary 2 to Proposition 7).

Let E be an A-algebra and \((a_{i})_{i\in I}\) a generating system of the A-module E (for example a basis). For E to be associative, it is necessary and sufficient that the \(a_{i}\) satisfy the associativity relations

\[
(a _ {i} a _ {j}) a _ {k} = a _ {i} (a _ {j} a _ {k}) \quad \text { for   all } i, j, k\tag{8}
\]

The mapping \((x, y, z) \mapsto (xy)z - x(yz)\) is an A-trilinear mapping

\[
\mathbf {E} \times \mathbf {E} \times \mathbf {E} \rightarrow \mathbf {E}
\]

and hence defines an A-linear mapping \(E \otimes_{A} E \otimes_{A} E \to E\) ; if the latter mapping is zero for all the elements \(a_{i} \otimes a_{j} \otimes a_{k}\) , which form a generating system of the A-module \(E \otimes_{A} E \otimes_{A} E\) , it is identically zero.

Similarly, for E to be commutative, it is necessary and sufficient that the \(a_{i}\) satisfy the commutativity relations

\[
a _ {i} a _ {j} = a _ {j} a _ {i} \quad \text { for   all } i, j.\tag{9}
\]

The proof is analogous, this time considering the A-bilinear mapping \((x,y)\mapsto xy - yx\) . Finally, for an element \(e\in E\) to be a unit element, it is necessary and sufficient that the \(a_{i}\) satisfy the relations

\[
a _ {i} = e a _ {i} = a _ {i} e \quad \text { for   all } i,\tag{10}
\]

as is seen this time by considering the A-linear mappings \(x \mapsto x - xe\) and \(x \mapsto x - ex\) .

When \((a_{i})_{i\in I}\) is a basis of E and \((\gamma_{ij}^{k})\) the corresponding family of constants of structure, relations (8) are equivalent to the relations

\[
\sum_ {r} \gamma_ {i j} ^ {r} \gamma_ {r k} ^ {s} = \sum_ {r} \gamma_ {i r} ^ {s} \gamma_ {j k} ^ {r}
\]

for all \(i, j, k, s\) . Similarly relations (9) are equivalent to \(\gamma_{ij}^{k} = \gamma_{ji}^{k}\) for all \(i, j, k\) .

Let \((a_{i})_{i\in I}\) be a basis of the A-algebra E; if \(\rho :A\to B\) is a ring homomorphism, \((a_{i}\otimes 1)\) is a basis of the B-algebra \(\rho^{*}(\mathbf{E}) = \mathbf{E}\otimes_{\mathbf{A}}\mathbf{B}\) (II, §5, no. 1, Proposition 4). If \((\gamma_{ij}^{k})\) is the family of constants of structure relative to the basis \((a_{i})\) , the family \((\rho (\gamma_{ij}^{k}))\) is the family of constants of structure of \(\rho^{*}(\mathbf{E})\) relative to the basis \((a_{i}\otimes 1)\) .

